\documentclass[11pt,a4paper]{article}
\usepackage[T1]{fontenc}
\usepackage{lmodern,amsmath,amssymb,amsthm,mathtools,microtype}
\usepackage[margin=27mm]{geometry}
\usepackage[colorlinks=true,linkcolor=blue,citecolor=blue,urlcolor=blue]{hyperref}
\setlength{\emergencystretch}{2em}
\newtheorem{theorem}{Theorem}[section]
\newtheorem{lemma}[theorem]{Lemma}
\newtheorem{corollary}[theorem]{Corollary}
\newcommand{\Q}{\mathbb Q}\newcommand{\Z}{\mathbb Z}
\newcommand{\Pp}{\mathbb P}
\DeclareMathOperator{\Jac}{Jac}
\DeclareMathOperator{\Hom}{Hom}
\DeclareMathOperator{\End}{End}
\title{Rational sections of fixed-input pole fibrations}
\author{}\date{16 November 2026}
\begin{document}\maketitle
\begin{abstract}
We determine when a fixed-input normalized pole equation has a rational
section in its conic parameter. For nonzero rational coefficients $C,D$,
an affine section exists exactly when $C$ or $-3D$ is a rational square.
When $C$ is nonsquare, all projective sections come from fixed points
of an explicit genus-one involution. The proof identifies the relevant
Jacobians as quadratic twists of a non-CM elliptic curve and accounts
for the possible lack of a rational origin on the second genus-one
curve. For prepared detector inputs, a section exists exactly when
$17(17-b)$ is a rational square. This classifies rational functions
of the parameter; it does not obstruct witnesses in individual rational
fibers or a common rational parameter for several inputs.
\end{abstract}

\section{The section criterion}

Fix $C,D\in\Q^\times$, let $\tau$ be transcendental over $\Q$, and put
\[
 k=\Q(\tau),\qquad e=\frac{\tau^2-D}{\tau^2+3D}.
\]
Let $F_{C,D}$ be the smooth projective curve over $k$ with affine equation
\begin{equation}\label{eq:pole}
                 q^2=D(s^2-Ce)(s^2+3Ce).
\end{equation}
Its four geometric branch points in the $s$ line are distinct, so it
has genus one. A rational section of the associated fibration over
$\Pp^1_\tau$ means a point of $F_{C,D}(k)$. An \emph{affine section}
means a pair $s,q\in k$ satisfying \eqref{eq:pole}.
All uses of ``square'' below mean a square in $\Q^\times$.

\begin{theorem}\label{thm:sections}
The following assertions hold.
\begin{enumerate}
\item If $C$ is nonsquare, the affine sections are exactly
\begin{equation}\label{eq:fixedzero}
              s=0,\qquad q=\pm Ce\sqrt{-3D},
\end{equation}
when $-3D$ is square, and there are no affine sections otherwise.
The only additional projective sections are the two points at infinity,
present exactly when $D$ is square.
\item If $C=c^2$ with $c\in\Q^\times$, there is the affine section
\begin{equation}\label{eq:squaresection}
             s=c,\qquad q=-\frac{4CD\tau}{\tau^2+3D}.
\end{equation}
\end{enumerate}
Consequently an affine section exists if and only if $C$ or $-3D$
is square. A projective section exists if and only if at least one
of $C,D,-3D$ is square.
\end{theorem}

We prove the theorem using explicit genus-one curves and descent.

\section{The trivializing curve and its Jacobian}

Let $B_D$ and $E_{C,D}$ denote the smooth projective curves with equations
\begin{equation}\label{eq:curves}
 B_D:\ r^2=\frac{\tau^2-D}{\tau^2+3D},\qquad
 E_{C,D}:\ v^2=D(u^2-C)(u^2+3C).
\end{equation}
Both have genus one. On $B_D$, the substitution
$z=r(\tau^2+3D)$ gives
\[
                  z^2=(\tau^2-D)(\tau^2+3D).
\]
The two points over $\tau=\infty$ have $r=\pm1$ and are rational.
Choose the point with $r=1$ as origin. An explicit Weierstrass model is
\begin{equation}\label{eq:jacobian}
 J_D:\quad Y^2=X^3-4DX^2+16D^2X,
\end{equation}
with mutually inverse rational maps
\begin{equation}\label{eq:weierstrassmaps}
 \begin{gathered}
 X=2(z+\tau^2+D),\qquad Y=2\tau X,\\
 \tau=\frac{Y}{2X},\qquad z=\frac X2-\tau^2-D.
 \end{gathered}
\end{equation}
They extend to an isomorphism of the smooth projective curves and
send the chosen origin to the Weierstrass origin. Direct calculation gives
\begin{equation}\label{eq:invariants}
 c_4=-512D^2,\qquad \Delta=-196608D^6,\qquad j(J_D)=\frac{2048}{3}.
\end{equation}
The $j$-invariant of a complex elliptic curve with complex
multiplication is an algebraic integer
\cite[Corollary 21.13]{Sutherland}. Since $2048/3$ is not an integer,
\begin{equation}\label{eq:endomorphisms}
                       \End_{\overline\Q}(J_D)=\Z.
\end{equation}

Write
\[
 \sigma(\tau,r)=(\tau,-r),\qquad \iota(u,v)=(-u,v)
\]
for the involutions of the two curves. Their quotient curves are of
genus zero, and each involution has four geometric fixed points.
In particular $\iota$ induces $[-1]$ on $\Jac(E_{C,D})$: after
choosing any geometric fixed point as origin, it is the nonidentity
origin-preserving involution of an elliptic curve. This assertion
does not require a rational fixed point.

\begin{lemma}\label{lem:twist}
The Jacobian of $E_{C,D}$ is the quadratic twist of $J_D$ by $C$.
If $C$ is nonsquare, then
\[
           \Hom_{\Q}(\Jac(B_D),\Jac(E_{C,D}))=0,
\]
and every $\Q$-morphism $B_D\to E_{C,D}$ is constant.
\end{lemma}
\begin{proof}
Choose $c\in\overline\Q$ with $c^2=C$. The formulas
\begin{equation}\label{eq:twistmaps}
 \begin{gathered}
 \phi_c:B_D\longrightarrow E_{C,D},\qquad
 u=\frac c r,\qquad
 v=-\frac{4CD\tau}{\tau^2-D}=-\tau(u^2-C),\\
 \tau=-\frac{v}{u^2-C},\qquad r=\frac c u
 \end{gathered}
\end{equation}
are mutually inverse rational maps, hence an isomorphism of smooth
projective curves over $\Q(c)$. For example,
\[
 u^2-C=\frac{4CD}{\tau^2-D},\qquad
 u^2+3C=\frac{4C\tau^2}{\tau^2-D},
\]
which verifies the target equation. Changing $c$ to $-c$ replaces
$\phi_c$ by $\iota\circ\phi_c$. If $C$ is nonsquare, the
nontrivial automorphism of $\Q(c)/\Q$ therefore sends the induced
Jacobian isomorphism $\Phi_c$ to $-\Phi_c$. If $C$ is square,
$\phi_c$ is already defined over $\Q$. These are precisely the
descent data for the quadratic twist by $C$; its Weierstrass equation is
\[
                  Y^2=X^3-4CDX^2+16C^2D^2X.
\]

By \eqref{eq:endomorphisms}, every geometric homomorphism between
these Jacobians is $n\Phi_c$ for some $n\in\Z$. If $C$ is
nonsquare, Galois conjugation sends it to $-n\Phi_c$. Hence the
only homomorphism defined over $\Q$ is zero.

Finally suppose $f:B_D\to E_{C,D}$ were nonconstant and defined
over $\Q$. The rational origin $O\in B_D(\Q)$ would give a
rational point $f(O)$ on $E_{C,D}$. Using these points as origins,
$f$ would be a nonzero homomorphism of elliptic curves by
\cite[Chapter II, Section 1, Proposition 1.5]{Milne}. It would
therefore induce a nonzero $\Q$-homomorphism of their Jacobians,
a contradiction. This argument does not assume a rational origin
on $E_{C,D}$ before the map $f$ is given.
\end{proof}

\section{Descent of sections}

Over the quadratic extension $\Q(B_D)/\Q(\tau)$, the substitutions
\begin{equation}\label{eq:trivialization}
                    u=\frac{s}{r},\qquad v=\frac{q}{r^2}
\end{equation}
identify $F_{C,D}$ with the constant curve $E_{C,D}$. It follows
that $F_{C,D}(\Q(\tau))$ is in bijection with the morphisms
\begin{equation}\label{eq:equivariance}
          f:B_D\longrightarrow E_{C,D}\quad\hbox{over }\Q,
          \qquad f\circ\sigma=\iota\circ f.
\end{equation}
Indeed a point of the generic fiber gives a rational map from $B_D$,
which extends over its smooth projective model; the condition in
\eqref{eq:equivariance} is precisely invariance of $s=ru$ and
$q=r^2v$ under $\sigma$. Conversely that condition gives descent
to $\Q(\tau)$. This includes the points at infinity on the
smooth projective models.

\begin{proof}[Proof of Theorem~\ref{thm:sections}]
If $C$ is nonsquare, Lemma~\ref{lem:twist} makes every map in
\eqref{eq:equivariance} constant. Its value is a rational point of
$E_{C,D}$, since $B_D$ has a rational point. Equivariance says that
this value is fixed by $\iota$.

In the affine chart the fixed points have $u=0$, and their ordinates
satisfy $v^2=-3DC^2$. They are rational exactly when $-3D$ is
square, and \eqref{eq:trivialization} gives
\eqref{eq:fixedzero}. At infinity use the coordinates
$t=1/u$ and $w=v/u^2$. Their equation and involution are
\[
 w^2=D(1-Ct^2)(1+3Ct^2),\qquad (t,w)\longmapsto(-t,w).
\]
The two points at $t=0$ are rational exactly when $D$ is square.
They correspond to the two infinity points of $F_{C,D}$, whose
local ordinate is $q/s^2=\pm\sqrt D$. These are all fixed points,
so the list of sections for nonsquare $C$ is exhaustive.

If $C=c^2$ with $c\in\Q^\times$, then
\[
 (1-e)(1+3e)=\frac{16D\tau^2}{(\tau^2+3D)^2}.
\]
Substituting $s=c$ in \eqref{eq:pole} proves
\eqref{eq:squaresection}. The two existence criteria now follow.
\end{proof}

The same maps also give a geometric description of the full surface:
it is birational to
\[
                  (B_D\times E_{C,D})/\langle\sigma\times\iota\rangle,
 \qquad s=ru,\quad q=r^2v.
\]
After choosing geometric fixed origins, this is the diagonal negation
quotient of a product of elliptic curves. The section argument above
retains the actual involutions, so it does not discard their possible
translation terms relative to rational origins.

\section{Prepared inputs and scope}

For a prepared pole input in \cite{Guards}, write $b=da^2$ with
$d,a\in\Q^\times$ and set
\[
                   D=17d,\qquad C=17(17-b).
\]
The preparation includes $0<b<17$ and $b\in\Q_{17}^{\times2}$.
Thus $D>0$, $C>0$, and $v_{17}(D)$ is odd. In particular neither
$D$ nor $-3D$ is a rational square.

\begin{corollary}\label{cor:prepared}
For a fixed prepared input, the normalized pole fibration has a
rational section if and only if $17(17-b)$ is a rational square.
For a finite list of fixed prepared inputs, a simultaneous tuple
of affine rational-function sections in the same parameter $\tau$
exists if and only if every $17(17-b_i)$ is a rational square.
When these conditions hold, \eqref{eq:squaresection} gives the tuple
explicitly, and each ordinate is nonzero as a rational function.
\end{corollary}
\begin{proof}
Apply Theorem~\ref{thm:sections} to each input. Projection to each
factor proves necessity for a simultaneous tuple, and the explicit
sections prove sufficiency.
\end{proof}

The criterion concerns $\Q(\tau)$, with $\tau$ transcendental.
Failure of a section does not imply that any particular rational
fiber has no rational point. It also does not exclude one common
rational value of $\tau$ for several prepared inputs, rational
multisections, or constructions that vary the detector inputs.
The criterion alone therefore gives no reduction in the number of
variables in a defining polynomial.

\section{Exact verification}

The reproducible checker is
\begin{center}\small
\texttt{negative\_results/computation/rational\_sections\_v42/check.py}.
\end{center}
It verifies the quotient substitution, both inverse pairs of maps,
the Weierstrass invariants and twist, the involution identities,
the complete fixed-point equations, and the explicit sections.
The non-CM, morphism, and Galois-descent steps are proved above with
their stated references; they are not inferred from numerical tests.

\begin{thebibliography}{9}\raggedright
\bibitem{Sutherland}
A.~V.~Sutherland, \emph{The Hilbert class polynomial},
MIT 18.783 Elliptic Curves, Spring 2017, Lecture~21,
1 May 2017, Corollary~21.13, p.~6,
\url{https://math.mit.edu/classes/18.783/2017/LectureNotes21.pdf}.
\bibitem{Milne}
J.~S.~Milne, \emph{Elliptic Curves}, second edition,
World Scientific, 2021,
Chapter~II, Section~1, Proposition~1.5,
\url{https://jmilne.org/math/Books/EC2.pdf}.
\bibitem{Guards}
\emph{Fixed pole fibers and independent detector indices},
project paper, version~42,
\texttt{papers/pole\_parameters/}\allowbreak
\texttt{common\_parameter\_guard\_reduction\_20261116\_v42.tex}.
\end{thebibliography}
\end{document}
